\cleardoublepage
\phantomsection
\addcontentsline{toc}{section}{Заключение}
\section*{Заключение}

Цель исследования~--- разработка метода реконструкции визуальных стимулов по одновременным записям фМРТ и ЭЭГ, использующего пространственно-временные характеристики обеих модальностей.

\textbf{Соответствие задач результатам.}
\begin{enumerate}
    \item \textit{Систематизация существующих подходов.} Глава~\ref{sec1} содержит структурированный анализ предметной области по идеям и инженерным решениям. Выделены два пробела в существующих работах: отсутствие двумодальных архитектур фМРТ--ЭЭГ для декодирования визуальных стимулов и отсутствие формализованного учёта гемодинамической задержки при формировании обучающих триплетов.
    \item \textit{Формализация задачи.} Глава~\ref{sec2} вводит совместное пространство $\mathcal{X}_f \times \mathcal{X}_e$, обучающую выборку триплетов с поправкой на гемодинамическую задержку $\Delta t$ и целевое отображение $\mathbf{f}: \mathcal{X}_f \times \mathcal{X}_e \to \mathcal{Y}$, оцениваемое метрикой CLIP-Score.
    \item \textit{Разработка метода.} Глава~\ref{sec3} описывает четыре компонента метода: предобработку с учётом $\Delta t$, архитектуру мультимодального энкодера ($\mathcal{E}_f$, $\mathcal{E}_e$, $\mathcal{F}$), контрастивное обучение с замороженным CLIP и двухстадийную диффузионную генерацию (априорная модель в CLIP-пространстве + предобученная SDXL с IP-Adapter).
    \item \textit{Экспериментальная проверка.} Глава~\ref{sec4} приводит две группы экспериментов. На вспомогательном наборе данных~\cite{Berezutskaya2022} оценено значение $\Delta t \approx 5$~с, подтверждена инвариантность весов линейной модели относительно испытуемых и значимость видеопризнаков. На основном наборе данных~\cite{Telesford2023} двумодальная конфигурация с априорной диффузионной моделью достигла CLIP-Score $0{,}644$--$0{,}668$ и превзошла одномодальные аналоги.
\end{enumerate}

\textbf{Основные результаты.}
\begin{itemize}
    \item Формализована задача двумодального декодирования визуальных стимулов с явным учётом гемодинамической задержки BOLD-сигнала при формировании обучающих триплетов.
    \item Предложена и реализована архитектура мультимодального энкодера, обучаемая контрастивно относительно CLIP-пространства, в сочетании с двухстадийной диффузионной генерацией.
    \item Получено эмпирическое подтверждение комплементарности модальностей: CLIP-Score двумодальной конфигурации систематически превышает CLIP-Score одномодальных конфигураций на всех трёх видеостимулах.
    \item Получено эмпирическое подтверждение значения гемодинамической задержки $\Delta t \approx 5$~с на независимом наборе данных и инвариантности линейной модели относительно испытуемых.
\end{itemize}

\textbf{Элементы научной новизны.}
\begin{itemize}
    \item Двумодальная архитектура декодирования визуальных стимулов по одновременным сигналам фМРТ--ЭЭГ с обучаемым модулем объединения модальностей.
    \item Компонент препроцессинга, явно учитывающий гемодинамическую задержку при формировании обучающих триплетов.
    \item Систематическое ablation-сравнение шести конфигураций (две модальности $\times$ три модальностных режима), количественно подтверждающее преимущество совместного использования модальностей и пользу априорной диффузионной модели.
\end{itemize}

\textbf{Теоретическая значимость} результатов состоит в эмпирическом подтверждении взаимодополнительности пространственной информации фМРТ и временной информации ЭЭГ в задаче визуального декодирования, а также в формализации связи между гемодинамической задержкой BOLD-сигнала и качеством контрастивного выравнивания нейронных представлений с CLIP-пространством.

\textbf{Практическая значимость.} Разработанный метод применим в исследованиях интерфейсов мозг--компьютер для визуального декодирования, в инструментах семантической реконструкции зрительной информации и в анализе нейронных откликов на натуралистические видеостимулы. Использование открытых наборов данных и общедоступных предобученных компонентов (CLIP, SDXL, IP-Adapter) обеспечивает воспроизводимость и переносимость подхода.

\textbf{Ограничения работы.} Реконструкция через CLIP-пространство обеспечивает семантическую близость, но не гарантирует пиксельной точности. Экспериментальная проверка проведена на одном основном наборе данных~\cite{Telesford2023}; перенос на иные типы стимулов и условия регистрации требует дополнительной валидации. Метод обучен и оценён на данных здоровых испытуемых; применение к пациентам с патологиями зрительной системы остаётся за рамками работы.

\textbf{Направления дальнейших исследований.}
\begin{itemize}
    \item Расширение экспериментальной части на статичные изображения, абстрактные паттерны и иные типы стимулов; валидация на дополнительных наборах данных одновременной регистрации фМРТ--ЭЭГ.
    \item Введение дополнительных метрик качества реконструкции, чувствительных к пиксельному уровню (например, LPIPS, SSIM, FID), в дополнение к CLIP-Score.
    \item Переход к нелинейным моделям оценки гемодинамической задержки, способным учитывать индивидуальные особенности BOLD-отклика и зависимость $\Delta t$ от области коры.
    \item Расширение модуля объединения модальностей: трансформер с кросс-вниманием, динамическое взвешивание модальностей в зависимости от типа стимула.
    \item Валидация на пациентах с нарушениями зрительной системы как потенциальная клиническая основа.
\end{itemize}
